\section{Generel Theory}
In this section, we will present the type of forward algorithm used to compute the likelihood 
and the pseudo residuals used to evaluate the model fit.


\subsection{Forward Algorithm} 
The forward algorithm implemented is called recursive filtering 
\citep{moler2024statistical} as it constructs the current 
statedistribution $p(X_t \mid Y_{1:t})$ recursively from the previous 
state distribution $p(X_t \mid Y_{1:t-1})$.
\begin{algorithm}[H]
\caption{Recursive Filtering Forward Algorithm}
\label{alg:forward}
\begin{algorithmic}[1]
\Require Observations $\mathbf{y} = (y_1, \ldots, y_T)$,
         covariates $\mathbf{x} = (x_1, \ldots, x_T)$,
         initial state distribution $\mathbf{u}_0 \in \mathbb{R}^{1 \times N}$,
         HMM parameters $\theta$
\Ensure Predicted distributions $\{\mathbf{u}_t\}$,
        filtered distributions $\{\mathbf{u}_{t|t}\}$,
        step likelihoods $\{f_t\}$

\For{$t = 1, \ldots, T$}
    \State $\Gamma_t \leftarrow \textsc{TransitionMatrix}(\theta,\, t,\, y_t,\, x_t)$
    \Comment{$N \times N$ transition matrix}
    \State $\mathbf{u}_{t \mid t - 1} \leftarrow \mathbf{u}_{t-1 \mid t-1} \cdot \Gamma_t$
    \Comment{Predicted state distribution, $1 \times N$}
    \State $\mathbf{g}_t \leftarrow \textsc{Density}(\theta,\, t,\, y_t,\, x_t)$
    \Comment{Emission densities for all states, $1 \times N$}
    \State $f_t \leftarrow \displaystyle\sum_{i=1}^{N} [\mathbf{u}_t \odot \mathbf{g}_t]_i$
    \Comment{Marginal likelihood at time $t$}
    \State $f_t \leftarrow \max(f_t,\; 10^{-10})$
    \Comment{Clip to avoid division by zero}
    \State $\mathbf{u}_{t \mid t} \leftarrow \dfrac{\mathbf{u}_t \odot \mathbf{g}_t}{f_t}$
    \Comment{Filtered state distribution (normalized)}
\EndFor
\State \Return $\{\mathbf{u}_t\},\; \{\mathbf{u}_{t \mid t}\},\; \{f_t\}$
\end{algorithmic}
\end{algorithm}

\noindent The log-likelihood is obtained by summing the log of the step likelihoods:
\begin{equation}
    \mathcal{L}(\theta) = \sum_{t=1}^{T} \log f_t
    \label{eq:log_likelihood}
\end{equation}

where $\odot$ denotes element-wise multiplication of vectos. 
\newline

The log-likelihood is obtained by summing the log of the step likelihoods:
\begin{equation}
    \mathcal{L}(\theta) = \sum_{t=1}^{T} \log f_t
    \label{eq:log_likelihood}
\end{equation}

\subsection{Forecast Pseudo Residuals}
A method to access the goodness-of-fit of the HMM is to compute the 
pseudo residuals. For a introduction to pseudo residuals, see chapter 6.2 in \citep{zucchini2016hidden}. \newline
The Forecast Pseudo Residuals used is the standard normal inverse of 
the cummulative distribution function (CDF) of
the one-step-ahead predictive distribution. If the model is is good, 
the pseudo residuals should be approximately standard normally distributed. 


Briefly desribed, a forecast pseudo residual for a HMM is defined as
\begin{equation}
    z_{t+1} = \Phi^{-1} (F_Y(y_{t+1}\mid Y_{1:t}))
    \label{eq:pseudo_residual}
\end{equation}

where $\Phi^{-1}$ is the inverse CDF of the standard normal distribution
and $F_Y$ is the predictive CDF of $y_{t+1}$ given $Y_{1:t}$, defined as
\begin{equation}
    F_Y(y_{t+1}\mid Y_{1:t}) = \sum_{i=1}^{N} u_{t+1 \mid t}^i F_i(y_{t+1} \mid X_t = i)
    \label{eq:predictive_cdf}
\end{equation}
where $X_t$ is the hidden state. \newline
Because $z_{t+1}$ now should follow a standard normal distribution, 
we can use normality tests to assess the goodness-of-fit of the model. \newline
In this project, we use the QQ-plot to visually assess the normality of the pseudo residuals. 

\subsection{Test Statistics}
Other methods to assess the goodness of fit of the HMM model 
include the AIC \citep{wiki:aic} and BIC \citep{wiki:bic} test statistics. These are defined as
\begin{equation}
    \text{AIC} = 2k - 2\mathcal{L}
    \label{eq:aic}
\end{equation}
\begin{equation}
    \text{BIC} = k \log(n) - 2\mathcal{L}
    \label{eq:bic}
\end{equation}
where $k$ is the number of parameters in the model, $n$ is the number of observations and $\mathcal{L}$ is the log-likelihood of the model. \newline
The AIC and BIC test statistics are used to compare different models,
where a lower value initiates a better fit. \newline
However, a lot of the models implemented are nested models, so a Likelihood Ratio Test (LRT) \citep{wiki:lrt} can also be
used to compare the models. The LRT test statistic is defined as
\begin{equation}
    \text{LRT} = 2(\mathcal{L}_1 - \mathcal{L}_0)
    \label{eq:lrt}
\end{equation}
where $\mathcal{L}_1$ is the log-likelihood of the more complex
model and $\mathcal{L}_0$ is the log-likelihood of the simpler model. \newline
The LRT test statistic follows a chi-squared distribution with degrees of freedom equal 
to the difference in the number of parameters between the two models. \newline




